\documentclass[10pt,a4paper]{article}
\usepackage[T1]{fontenc}
\usepackage{lmodern,amsmath,amssymb,booktabs,array,xcolor,microtype}
\usepackage[margin=16mm,top=14mm,bottom=16mm]{geometry}
\usepackage{fancyhdr,tikz,graphicx}
\usepackage[hidelinks]{hyperref}
\definecolor{navy}{HTML}{17324D}
\definecolor{teal}{HTML}{176E78}
\definecolor{pale}{HTML}{EEF5F7}
\definecolor{muted}{HTML}{4B5963}
\definecolor{rulegray}{HTML}{CAD5DB}
\setlength{\parindent}{0pt}
\setlength{\parskip}{3pt}
\setlength{\abovedisplayskip}{4pt}
\setlength{\belowdisplayskip}{4pt}
\setlength{\abovedisplayshortskip}{3pt}
\setlength{\belowdisplayshortskip}{3pt}
\pagestyle{fancy}\fancyhf{}
\renewcommand{\headrulewidth}{0pt}
\fancyfoot[L]{\scriptsize\color{muted}QVDF-E \quad | \quad Arizona detector 141 worked example}
\fancyfoot[R]{\scriptsize\color{muted}\thepage\ / 7}
\newcommand{\method}[2]{\par\vspace{4pt}{\color{teal}\bfseries #1\quad #2}\par\vspace{1pt}}
\newcommand{\notebox}[1]{\par\vspace{4pt}\begingroup\setlength{\fboxsep}{7pt}\colorbox{pale}{\parbox{\dimexpr\linewidth-14pt\relax}{#1}}\endgroup\par}
\hypersetup{pdftitle={QVDF-E: Arizona detector 141 emission diagnostic},pdfauthor={Mohammad Abbasi},pdfsubject={Equations and verified average-weekday results}}
\definecolor{methodO}{HTML}{242B33}
\definecolor{methodS}{HTML}{2563EB}
\definecolor{methodR}{HTML}{7B3294}
\definecolor{methodM}{HTML}{168C69}
\definecolor{methodA}{HTML}{E69F00}\newcommand{\legendkey}[2]{{\scriptsize\textcolor{#1}{\rule[2pt]{10pt}{1.3pt}}\ #2}}
\begin{document}
\fontsize{9.3}{11.1}\selectfont
{\fontsize{19}{22}\selectfont\bfseries\color{navy}Five emission calculations}\par
{\color{muted}One whole episode, one pollutant. O, S, R, M and A identify the existing calculations.}

\textbf{Notation and units.} Bin $i$ has speed $v_i$ and count $d_i=q_i\Delta t$, with $\Delta t=1/12$ h.
$D=\sum_i d_i$; $L$ is link length; $C$ is hourly capacity; $T_f=L/v_f$.
$H$ denotes vehicle-hours, $P$ episode duration and $T$ per-vehicle travel time.
$r_0(v)$ is the zero-acceleration lookup rate in g/(veh\,h); $E$ is in grams. Counts and capacity use the same lane basis.

\textbf{Shared fitted quantities for R, M and A.}
\[
x=\frac{D}{C},\qquad \widehat P=\max(f_dx^n,x),\qquad
\widehat w_{\max}=T_f f_p\widehat P^{s},\qquad
\widehat{\bar w}=\theta\widehat w_{\max},\qquad
H_{\rm pred}=D(T_f+\widehat{\bar w}).
\]
$f_d,n,f_p,s,\theta$ are calibrated by directional corridor and period, with the existing sparse-period fallback.

\method{O}{Recorded-speed reference}
Assign each bin's counted vehicles the complete-link travel time $L/v_i$, then sum:
\[
H_i=d_i\frac{L}{v_i},\qquad H_{\rm obs}=\sum_i H_i,\qquad
\boxed{E_O=\sum_i H_i r_0(v_i).}
\]
The five-minute interval supplies the count; emission exposure is $L/v_i$, not five minutes.

\method{S}{Episode mean-speed estimate}
Use the count-weighted harmonic mean, preserving the reference's total vehicle-hours:
\[
\bar v_{\rm obs}=\frac{D}{\sum_i d_i/v_i}=\frac{DL}{H_{\rm obs}},\qquad
\boxed{E_S=H_{\rm obs}\,r_0(\bar v_{\rm obs}).}
\]

\method{R}{Reconstructed whole-link speeds}
For normalized entry time $u\in[0,1]$, define $w(u)=\widehat w_{\max}h(u)$ and
$v_R(u)=L/[T_f+w(u)]$. The symmetric profile is
\[
h(u)=\begin{cases}
[4u(1-u)]^a,&\theta\le 2/3,\quad \int_0^1h(u)\,du=\theta\ \text{determines }a,\\
1-|2u-1|^b,&\theta>2/3,\quad b=\theta/(1-\theta).
\end{cases}
\]
The entering-vehicle weight is $dN=D[1+(\widehat w_{\max}/\widehat P)h'(u)]\,du$. Thus
\[
\boxed{E_R=\int_0^1[T_f+w(u)]\,r_0\!\left(\frac{L}{T_f+w(u)}\right)\,dN.}
\]
The integral uses 384 quadrature points; $\int dN=D$ and $\int[T_f+w(u)]dN=H_{\rm pred}$.
These are trip-average speeds by entry time, not individual operating-speed trajectories.

\method{M}{Two-speed lookup}
Infer the queue speed from discharge and the triangular fundamental diagram:
\[
\widehat\mu=\frac{D}{\widehat P},\qquad
\omega=\frac{C}{k_j-C/v_f},\qquad
\widehat v_q=\frac{\widehat\mu}{k_j-\widehat\mu/\omega},\qquad k_j=220\ \text{veh/mi/lane}.
\]
Allocate mean travel time between two operating speeds while preserving distance:
\[
\bar T_q=\frac{v_f\widehat{\bar w}}{v_f-\widehat v_q},\qquad
\bar T_{\rm free}=T_f-\frac{\widehat v_q\widehat{\bar w}}{v_f-\widehat v_q},\qquad
\boxed{E_M=D[\bar T_{\rm free}r_0(v_f)+\bar T_qr_0(\widehat v_q)].}
\]
Here $\bar T_{\rm free}+\bar T_q=T_f+\widehat{\bar w}$ and $v_f\bar T_{\rm free}+\widehat v_q\bar T_q=L$.

\method{A}{Two-speed cubic}
Keep M's speeds and times; replace the lookup rate by its fitted cubic:
\[
r_c(v)=c_0+c_1v+c_2v^2+c_3v^3,\qquad
\boxed{E_A=D[\bar T_{\rm free}r_c(v_f)+\bar T_qr_c(\widehat v_q)].}
\]

\notebox{\textbf{Evaluation: episode-total WAPE.} For method $m\in\{S,R,M,A\}$ and the same $N=137$ episodes,
\[
\boxed{\operatorname{WAPE}_m=100\frac{\sum_{k=1}^{N}|E_{m,k}-E_{O,k}|}{\sum_{k=1}^{N}E_{O,k}}.}
\]
Absolute errors are calculated before pooling episodes. O is a detector-speed-based calculation,
not measured emissions or observed vehicle trajectories. All five calculations omit acceleration effects.}
